\chapter{Proposal Feedback and Iteration Plan}
\label{ch:iteration-plan}

\section{Proposal Feedback Summary}
Supervisor feedback was received over five meetings with Dr.~Suleman Khan. In the first meeting the team was
asked to review papers related to the project idea; in the second, the project scope and modules were revised;
in the third, dataset augmentation and preprocessing were discussed and a dashboard showing live updates was
requested; and in the fifth, the supervisor asked for multiple real Pakistani traffic scenarios in the
simulation. The meeting-wise feedback and the action taken are listed in the appendix. The action register below
(Table~\ref{tab:action-register}) lists the research- and scope-driven actions taken from the proposal, with
their status after Iteration~2.

\section{Proposal Action Register}
\begin{longtable}{p{0.05\textwidth}p{0.27\textwidth}p{0.29\textwidth}p{0.29\textwidth}}
\caption{Proposal action register.}\label{tab:action-register}\\
\toprule
No. & Required action / issue & Action taken & Status / evidence\\
\midrule
\endfirsthead
\toprule
No. & Required action / issue & Action taken & Status / evidence\\
\midrule
\endhead
1 & Define a focused contribution beyond generic adaptive traffic lights. &
Focus on ambulance priority with fairness-aware recovery. & Done; scope in Chapter~1.\\
2 & Justify the research gap with a targeted comparison. &
Compared fairness-control, emergency-preemption and detection literature. & Done; Chapter~2.\\
3 & Do not claim detection without labelled data. &
Trained a labelled ambulance / other-vehicle classifier; checked the splits for duplicates and evaluated it
once on the held-out test split. & Done; Section~\ref{sec:it2-classifier}. Dataset-licence notes pending.\\
4 & Evaluate impact on ordinary traffic as well as ambulance travel time. &
1\,000 held-out SUMO runs comparing fixed-time, actuated, fixed-time with priority and EquiFlow on car wait,
tail wait, worst approach and Jain index. & Done; Section~\ref{sec:it2-results}.\\
5 & Make recovery after preemption an explicit feature. &
Queue-sized recovery greens (longest wait first) and a fairness-aware MARL base controller, both implemented
and ablated. & Implemented; no measurable fairness benefit yet (Section~\ref{sec:it2-ablation}).\\
\bottomrule
\end{longtable}

\section{Iteration Strategy}
We develop the system incrementally. Each iteration produces a runnable artifact and supporting evidence.
Iteration~1 established the foundations that everything else depends on: sensing (ambulance recognition and a
dashboard view) and simulation (SUMO scenarios and baselines). Iteration~2 adds tracking and control and
integrates the modules. A learning-based controller is introduced only on top of the stable simulation state,
action interface and baseline metrics from Iteration~1.

\section{Iteration 1 Plan}
\textbf{Objectives:} build the sensing and simulation foundations and a modular repository.
\begin{itemize}
\item Train a YOLO26 ambulance recognition model on a labelled dataset. \emph{Done.}
\item Display detections on video in a live dashboard. \emph{Done.}
\item Build SUMO scenarios (Islamabad intersection types, 5 vehicle classes, 4 demand scenarios). \emph{Done.}
\item Implement fixed-time and actuated baseline controllers. \emph{Done.}
\item Organise the repository as vision, marl, sumo, comm, dashboard and docs. \emph{Done.}
\end{itemize}
\textbf{Deliverables and evidence:} trained model and training logs, dashboard demonstration, SUMO toolkit,
repository (Chapter~\ref{ch:iteration2} builds on these).

\section{Iteration 2 Plan}
\textbf{Objectives:} add tracking and control, then integrate the modules.
\begin{itemize}
\item Add ByteTrack tracking so detections persist across frames.
\item Implement route-aware preemption with safe phase transitions and a green corridor where feasible.
\item Implement and train a fairness-aware MARL controller and a recovery policy.
\item Connect the modules through FastAPI and MQTT and link the dashboard to the live system.
\item Complete the formal detector evaluation (test split, latency) and run repeatable baseline comparisons.
\end{itemize}
\textbf{Deliverables:} integrated prototype, experiment scripts, preliminary results and documented failure
cases. \textbf{Exit criterion:} a scripted ambulance scenario triggers and completes a logged preemption with
recovery. The outcome of each item is reported in Chapter~\ref{ch:iteration2}.

\section{Work Breakdown and Responsibilities}
\begin{table}[htbp]
\centering
\small
\caption{Work breakdown and responsibilities.}
\label{tab:work-breakdown}
\begin{tabular}{p{0.24\textwidth}p{0.36\textwidth}p{0.3\textwidth}}
\toprule
Workstream & Primary responsibility & Supporting role / evidence\\
\midrule
Research and requirements & Literature review, requirements, citations. & Literature matrix and meeting notes.\\
Data and detection (vision) & Dataset, model training, detector evaluation, tracking. &
Training logs, split manifest, test-split evaluation, tracking scores.\\
Simulation and control (sumo, marl) & SUMO networks, demand, baselines, emergency routes, MARL, priority and
recovery. & Configuration, run manifest and reproducible runs.\\
Integration (comm, dashboard) & FastAPI/MQTT interfaces, dashboard, end-to-end tests. &
Code, event logs, scripted scenario, demonstrations.\\
Documentation and release (docs) & Report, version control, test evidence, packaging. &
Repository, pinned versions, results documents.\\
\bottomrule
\end{tabular}
\end{table}

\section{Progress Against Plan}
\begin{longtable}{p{0.27\textwidth}p{0.16\textwidth}p{0.47\textwidth}}
\caption{Progress against plan after Iteration 2.}\label{tab:progress}\\
\toprule
Work item & Status & Evidence\\
\midrule
\endfirsthead
\toprule
Work item & Status & Evidence\\
\midrule
\endhead
Problem and research gap & Done & Literature review; Chapter~2.\\
Dataset preparation & Done & 32\,593 labelled crops in train/val/test after removing 39 duplicate or
mislabelled files; split manifest with leakage check.\\
YOLO26 ambulance recognition & Done & Mid-report model cls\_v4-2; retrained cls\_v5\_img128 now used
(Table~\ref{tab:it2-cls-val}).\\
Live detection dashboard & Done & Next.js dashboard linked to the live backend.\\
SUMO toolkit and baselines & Done & Fixed-time and actuated controllers; Islamabad junction networks.\\
Repository structure & Done & vision, marl, sumo, comm, dashboard, backend, docs.\\
ByteTrack tracking & Done & 12 real videos scored; 7 of 7 ambulance passes confirmed
(Table~\ref{tab:it2-tracking}).\\
Route-aware preemption & Done & Safe yellow and all-red, green corridor on the multi-junction networks.\\
Fair recovery / MARL & Done, no measurable benefit & Ablations show no significant effect of recovery;
MARL ties the heuristic (Section~\ref{sec:it2-ablation}).\\
FastAPI / MQTT integration & Done & REST, WebSocket, MQTT, CSV event log; scripted scenario passes.\\
Test-split and latency evaluation & Done & 93.5\% recall, 96.5\% precision on the test split
(Table~\ref{tab:it2-cls-test}); about 45\,ms per frame on an RTX 4050 (Table~\ref{tab:it2-latency}).\\
Repeatable baseline comparisons & Done & 1\,000 one-hour runs on held-out seeds with 95\% confidence intervals
and paired differences.\\
\bottomrule
\end{longtable}

\section{Delays, Deviations and Recovery Plan}
Iteration~1 was scoped to the two stages that other stages depend on (sensing and simulation), so tracking,
preemption, MARL and integration were deliberately moved to Iteration~2 instead of being started in parallel,
together with the formal detector evaluation.

In Iteration~2 every planned item was delivered, with these deviations:
\begin{itemize}
\item \textbf{Detector evaluation came last.} It needed the full crops dataset and a GPU, so it ran after the
control and simulation work. The duplicate check it required found four mislabelled validation crops, which were
removed; the classifier was then retrained and the better model (cls\_v5\_img128) was tested once on the held-out
split.
\item \textbf{The confirmation vote was re-tuned.} With the new classifier the original vote (5 of 8 frames)
missed two of the seven ambulance passes in the real videos; an offline sweep moved it to 3 of 8. Because the same
videos chose and scored the setting, a leave-one-video-out estimate is reported alongside.
\item \textbf{One labelled video (test7) has no source file} and was dropped from the tracking evaluation.
\item \textbf{Fair recovery and MARL did not show the expected benefit.} Both are implemented and evaluated, but
recovery has no statistically significant effect and MARL only ties the hand-tuned adaptive controller.
\end{itemize}
No dated project log was kept, so schedule slippage is not quantified. The recovery plan for FYP-2 follows the
risks identified at the mid-report: the working rule-based controller and the SUMO evaluation remain the
baseline, the heuristic recovery stays in place while a recovery policy with a measurable fairness benefit is
studied, MARL training is extended, and more labelled real video and local ambulance images are collected to
make the detection results statistically meaningful.
